% !TeX root = ../教材题库_全部条目_转写工作稿.tex
% Source: XJ-B1-C04-HW-P104-Y0598
% 第4章 / 4.1 实数指数幂和幂函数 / 教材 p.104 / PDF 第111页
% OCR未校对：公式与文字均可能有误；未完成数学TeX转写；禁止直接发给学生
% 原题及相邻原解截图（核对用）：../题目原图/XJ-B1-C04-HW-P104-Y0598-01.webp
% 本文件只在转写完成后作为题干使用；OCR 条目尚不是可靠题干。
\begin{bankitem}{习题4.1 第2题}{第4章 / 4.1 实数指数幂和幂函数 / 教材 p.104 / PDF 第111页}{XJ-B1-C04-HW-P104-Y0598}{OCR未校对：公式与文字均可能有误；未完成数学TeX转写；禁止直接发给学生}
\begin{ocrdraft}
[教材p.104 / PDF第111页]
\par
2．用分数指数幂的形式表示下列各式的化简结果（其中a，b均为正数）：
\par
2
\par
（2）
\par
（3）2a3乃
\end{ocrdraft}
\end{bankitem}
